\section{Problem Formulation}
\label{sec:problem-formulation}

\paragraph{Open-ended Web exploration.}
Let $\mathcal{W}$ denote a Web application whose internal business state and
transition rules are unknown. At interaction step $t$, an agent receives a GUI
observation $o_t \in \mathcal{O}$ and selects a browser action
$a_t \in \mathcal{A}$. Executing $a_t$ produces a new observation $o_{t+1}$
and executor-reported status $x_t$. The resulting interaction record is
$\tau_t=(o_t,a_t,x_t,o_{t+1})$. Unlike task-conditioned Web agents, the agent
is not given a user goal to complete. It instead explores the application to
induce a reusable functional model from interaction traces. The model is
limited to states and outcomes that are observable through the GUI or saved
interaction artifacts; it is not intended to recover hidden application logic.

\paragraph{Testable functional hypotheses.}
From an observation, the agent may propose a functional hypothesis
\begin{equation}
    h_i=(\ell_i,f_i,\hat{y}_i),
\end{equation}
where $\ell_i$ is a semantic location, $f_i$ is a high-level function,
and $\hat{y}_i$ is an expected observable outcome. The claim states that
executing $f_i$ at $\ell_i$ is expected to produce $\hat{y}_i$. Semantic
locations and high-level functions describe application-level concepts rather
than individual DOM elements. The model may additionally retain applicability
metadata, including location constraints and observed direct action
dependencies. Such dependencies record interaction-supported ordering
relationships; they are not part of the outcome claim itself and are not
asserted to be necessary or sufficient business preconditions.

The expected outcome $\hat{y}_i$ is generated together with the hypothesis and
frozen before execution. It must identify a change that can be checked in a
subsequent GUI observation or saved observable state. Freezing $\hat{y}_i$
prevents the claimed outcome from being retrofitted after observing the
execution result. A proposed hypothesis is a candidate for verification, not
yet a fact about the application.

\paragraph{Execution evidence and functional outcome.}
For an attempted hypothesis $h_i$, the system stores an evidence packet
\begin{equation}
    e_i=(o_i^{\mathrm{pre}}, a_i, x_i,
         o_i^{\mathrm{post}}, m_i),
\end{equation}
where $o_i^{\mathrm{pre}}$ and $o_i^{\mathrm{post}}$ are the before and after
observations, $a_i$ is the executed interaction, $x_i$ is the executor-reported
status, and $m_i$ contains available metadata and evidence references. An
evidence judgment
\begin{equation}
    J(h_i,e_i) \in \{\textsc{Supported},\textsc{Unsupported},
    \textsc{Unresolved}\}
\end{equation}
evaluates whether $e_i$ supports the frozen expected outcome $\hat{y}_i$.
\textsc{Unsupported} requires sufficiently complete evidence that does not
support the expected outcome. A hypothesis that has not been attempted, whose
execution failed before producing diagnostic evidence, or whose evidence is
missing or only partial remains \textsc{Unresolved}. We separately record
whether an executed interaction preserves the current semantic location or
transitions to another one.

This definition distinguishes three claims that are otherwise easy to
conflate: proposing $h_i$ establishes only that a candidate was discovered;
an executor success $x_i$ establishes only that the interaction was reported
as completed; and $J(h_i,e_i)=\textsc{Supported}$ establishes that the available
before--after evidence supports the expected functional outcome. Neither
candidate discovery nor executor-reported success alone is sufficient for
functional knowledge admission.

\paragraph{Evidence judgment and knowledge admission.}
The three-way judgment is distinct from the binary admission indicator
$A(h_i)\in\{0,1\}$. A hypothesis is admitted to the downstream-usable
functional knowledge set only when its evidence is complete and supports its
frozen expected outcome:
\begin{equation}
    A(h_i)=1
    \quad\Longleftrightarrow\quad
    \operatorname{complete}(e_i) \land
    J(h_i,e_i)=\textsc{Supported}.
    \label{eq:admission}
\end{equation}
When $A(h_i)=0$, the hypothesis is not admitted under the current evidence.
This does not by itself mean that the proposed function is false. Non-admitted
records and diagnostic reasons are retained for inspection and possible
continuation; in particular, an unresolved attempt is not treated as negative
evidence about the proposed function. We call the set of hypotheses, their
evidence judgments and admission indicators, and their linked interaction
evidence the evidence-grounded functional model $\mathcal{M}$.

\paragraph{Context-conditioned interaction risk.}
Before executing a selected browser action $a_t$ under observation $o_t$, the
system assigns an interaction-level risk annotation
\begin{equation}
    \rho_t=(b_t,k_t,q_t)=R(o_t,a_t,\mathcal{K}),
\end{equation}
where $b_t$ is a binary potential-risk judgment, $k_t$ is one primary risk type
(or null when no risk is identified), $q_t$ is supporting interface evidence,
and $\mathcal{K}$ is a versioned risk taxonomy. The object of this judgment is
whether executing this selected interaction in the current GUI context may
perform, enter, or prepare a risk-sensitive operation. The implemented judge
receives the pre-action screenshot and a semantic label for $a_t$. It does not
assign an intrinsic, context-free risk property to $f_i$, assess an entire
multi-action verification attempt as one unit, or enumerate every possible
action visible on the screen.

The annotation $\rho_t$ is attached to the interaction record $\tau_t$. When
$a_t$ is selected to test $h_i$, the record also links $\rho_t$ to that
functional claim and its subsequent outcome evidence. Replay actions are
assessed individually and linked to their replay steps, even when they do not
propose a new claim. In the system evaluated here, risk assessment runs in
shadow mode: it records a pre-execution judgment but does not block or modify
the selected action. Consequently, the formulation supports risk awareness and
auditability, but does not by itself imply safer execution or an end-to-end
safety guarantee.

The overall problem is therefore to induce $\mathcal{M}$ through open-ended
interaction while (i) admitting only functional claims supported by observable
execution evidence, (ii) continuing to produce useful supported knowledge
under a bounded exploration budget, and (iii) associating each selected
evidence-gathering action with a context-conditioned pre-execution risk record.
